\documentclass[11pt,a4paper]{report}

% ===================== PACKAGES =====================
\usepackage[utf8]{inputenc}
\usepackage[T1]{fontenc}
\usepackage[french]{babel}
\usepackage{amsmath,amssymb,amsthm}
\usepackage{graphicx}
\usepackage{geometry}
\geometry{margin=2.5cm}
\usepackage{hyperref}
\hypersetup{colorlinks=true, linkcolor=black, citecolor=black, urlcolor=blue}
\usepackage{booktabs}
\usepackage{caption}
\usepackage{listings}
\usepackage{xcolor}
\usepackage{fancyhdr}
\usepackage[numbers,sort&compress]{natbib}

% ===================== MISE EN PAGE =====================
\pagestyle{fancy}
\fancyhf{}
\fancyhead[L]{\leftmark}
\fancyfoot[C]{\thepage}

\lstset{
    language=Python,
    basicstyle=\ttfamily\small,
    keywordstyle=\color{blue},
    commentstyle=\color{gray},
    stringstyle=\color{red},
    numbers=left,
    numberstyle=\tiny,
    frame=single,
    breaklines=true
}

% ===================== INFOS DOCUMENT =====================
\title{
    \Large Université de Strasbourg \\
    \large Master Mathématiques et Applications -- Parcours Statistique \\
    \vspace{1cm}
    \Huge \textbf{Identification individuelle et analyse du réseau social \\ de macaques japonais par apprentissage profond} \\
    \vspace{0.3cm}
    \large Projet AI-SARU -- Population de Koshima
}
\author{Prénom NOM}
\date{Année universitaire 2025--2026}

% ===================== DEBUT DOCUMENT =====================
\begin{document}

\maketitle
\thispagestyle{empty}

\chapter*{Remerciements}
\thispagestyle{empty}
% TODO : rédiger vos remerciements

\newpage

\tableofcontents
\thispagestyle{empty}

\newpage


% =====================================================================
% INTRODUCTION
% =====================================================================
\chapter*{Introduction}
\addcontentsline{toc}{chapter}{Introduction}

\section*{Contexte scientifique et présentation du projet}
% TODO : présenter le laboratoire, ses thématiques de recherche, et le projet
% AI-SARU (pipeline en 3 modules : détection des comportements, reconnaissance
% des individus, mesure de proximité et réseau)

\section*{Contexte du stage et problématique}
% TODO : histoire ayant mené à ce sujet ; problématique -- accompagner un
% chercheur en éthologie dans la réussite de son modèle de deep learning (YOLO)
% pour identifier individuellement des macaques et en déduire un réseau social

\section*{Objectifs du stage}
% TODO : objectifs opérationnels -- annotation/correction des données de
% Koshima, ré-entraînement du modèle, construction et évaluation de la matrice
% de proximité sociale

\section*{Organisation du rapport}
% TODO : brève annonce du plan (chapitres 1 à 5)


% =====================================================================
% CHAPITRE 1 — MÉTHODES DE DÉTECTION AUTOMATIQUE PAR DEEP LEARNING
% =====================================================================
\chapter{Méthodes de détection automatique par deep learning}
\label{chap:methode}

\section{Introduction au machine learning et au deep learning}

\paragraph{Apprentissage supervisé.}

On se place dans le cadre de l'apprentissage supervisé : on dispose d'un ensemble
d'exemples $(X_i, Y_i)_{i=1}^{n}$, où $X_i$ est une image et $Y_i \in \{1,\dots,K\}$
l'identité connue de l'individu photographié parmi les $K$ individus de la population
étudiée. L'objectif est d'estimer une fonction $\hat{f}_\theta$, paramétrée par $\theta$,
qui approche au mieux la relation entre $X$ et $Y$. Il s'agit d'un problème de
\emph{classification}, la variable $Y$ étant catégorielle (à distinguer d'un problème
de régression, où $Y$ serait continue).

\paragraph{Modèle de classification et fonction de perte.}

Pour une image $X_i$, le modèle produit un vecteur de probabilités
\begin{equation}
    \hat{p}(X_i) = (\hat{p}_1, \dots, \hat{p}_K), \qquad \sum_{k=1}^{K} \hat{p}_k = 1,
\end{equation}
et l'identité prédite est $\hat{Y}_i = \arg\max_k \hat{p}_k$.

L'apprentissage consiste à ajuster $\theta$ en minimisant une fonction de perte,
l'entropie croisée :
\begin{equation}
    \mathcal{L}(\theta) = -\frac{1}{n} \sum_{i=1}^{n} \log \hat{p}_{\theta, Y_i}(X_i).
\end{equation}

Ce choix n'est pas arbitraire : sous l'hypothèse d'observations indépendantes, la
vraisemblance du modèle s'écrit
\begin{equation}
    \mathcal{V}(\theta) = \prod_{i=1}^{n} \hat{p}_{\theta, Y_i}(X_i),
\end{equation}
et minimiser $\mathcal{L}(\theta)$ revient exactement à maximiser la log-vraisemblance
$\log \mathcal{V}(\theta)$, à un facteur $-1/n$ près. L'entraînement du réseau de
neurones s'interprète donc comme une \textbf{estimation par maximum de vraisemblance}.

\paragraph{L'opération de convolution.}

La fonction $\hat{f}_\theta$ est ici réalisée par un \textbf{réseau de neurones
convolutif} (CNN). L'opération de base de ce type de réseau est la
\textbf{convolution discrète} : un petit filtre (ou noyau) $W \in \mathbb{R}^{k \times k}$,
de paramètres appartenant à $\theta$, glisse sur l'image d'entrée $X$ et produit une
carte de caractéristiques $F$ définie, pour chaque position $(i,j)$, par
\begin{equation}
    F(i,j) = \sum_{u=1}^{k} \sum_{v=1}^{k} W(u,v) \, X(i+u, j+v) \; + \; b,
\end{equation}
où $b$ est un biais. Ce résultat est ensuite transformé par une fonction d'activation
non linéaire $\phi$ (typiquement la \emph{ReLU}, $\phi(x) = \max(0,x)$), donnant une
carte activée $\phi(F)$. Un réseau convolutif empile ainsi plusieurs couches de ce
type : les premières couches détectent des motifs simples (contours, textures), les
couches plus profondes combinent ces motifs en des représentations de plus haut
niveau (formes de visage, motifs propres à un individu). L'ensemble des poids $W$ et
biais $b$ de toutes les couches constitue le vecteur de paramètres $\theta$ à
estimer.

\paragraph{Estimation des paramètres par descente de gradient.}

L'estimation de $\hat{\theta} = \arg\min_\theta \mathcal{L}(\theta)$ ne peut pas se
faire par une formule fermée, la fonction $\mathcal{L}$ étant hautement non convexe
et de très grande dimension (typiquement plusieurs millions de paramètres). On
procède donc par \textbf{descente de gradient} : partant d'une valeur initiale
$\theta^{(0)}$, on met à jour itérativement les paramètres dans la direction opposée
au gradient de la perte,
\begin{equation}
    \theta^{(t+1)} = \theta^{(t)} - \eta \, \nabla_\theta \mathcal{L}(\theta^{(t)}),
\end{equation}
où $\eta > 0$ est le taux d'apprentissage (\emph{learning rate}). En pratique, le
gradient n'est pas calculé sur l'ensemble des $n$ exemples à chaque itération (ce
qui serait trop coûteux), mais sur un sous-échantillon aléatoire (un \emph{batch}) :
c'est la \textbf{descente de gradient stochastique} (SGD), ou l'une de ses variantes
(Adam, RMSProp), utilisée en pratique par la quasi-totalité des réseaux de neurones,
y compris YOLO. Le calcul du gradient $\nabla_\theta \mathcal{L}$, malgré la
profondeur du réseau, est rendu possible par l'algorithme de
\textbf{rétropropagation} (\emph{backpropagation}), qui applique la règle de
dérivation en chaîne couche par couche, depuis la sortie du réseau jusqu'à ses
premiers paramètres.

Le \textbf{ré-entraînement} (\emph{fine-tuning}) réalisé pendant le stage consiste
précisément à reprendre cette procédure d'optimisation, non pas depuis une
initialisation aléatoire $\theta^{(0)}$, mais depuis les paramètres déjà appris sur
la population de Shodoshima, en poursuivant la descente de gradient sur les
nouvelles données annotées de Koshima. Cette approche exploite le fait que les
premières couches du réseau (détection de motifs visuels génériques) restent
largement valides d'une population à l'autre, seules les couches les plus profondes
nécessitant un réajustement substantiel pour s'adapter aux nouveaux individus.

\section{Principe de la détection d'objets et du modèle YOLO}

Le modèle utilisé pour la reconnaissance faciale des macaques est un modèle
\textbf{YOLO} (\emph{You Only Look Once}) \citep{redmon2016}. Contrairement à un
simple classifieur qui recevrait en entrée une image déjà recadrée sur un visage,
YOLO est un modèle de \textbf{détection d'objets} : à partir d'une image complète, il
prédit simultanément, en une seule passe du réseau,
\begin{itemize}
    \item la \textbf{localisation} des objets présents (ici, les coordonnées des
    boîtes englobantes autour des visages détectés) ;
    \item la \textbf{classification} de chaque objet détecté, c'est-à-dire
    l'identité de l'individu parmi les $K$ individus connus.
\end{itemize}
C'est ce mécanisme qui explique le fonctionnement observé dans l'interface
d'annotation : pour chaque visage localisé, le modèle renvoie les identités
candidates classées par score de confiance, ce score correspondant à la probabilité
de classe $\hat{p}_k$ associée à la détection.

\paragraph{Découpage en grille.}

Le principe central de YOLO consiste à découper l'image d'entrée en une grille de
$S \times S$ cellules. Chaque cellule est responsable de détecter un objet si le
centre de cet objet tombe dans son emprise. L'ensemble de l'image est ainsi traité
en une seule passe du réseau, contrairement aux approches plus anciennes qui
balayaient l'image par fenêtres glissantes successives -- d'où le nom de la méthode.

\paragraph{Mesure de qualité d'une détection : l'IoU.}

Pour évaluer si une boîte englobante prédite correspond bien à un objet réel, on
utilise l'\textbf{intersection sur union} (\emph{IoU}) entre la boîte prédite et la
boîte de vérité terrain :
\begin{equation}
    \text{IoU} = \frac{\text{aire}(\text{boîte prédite} \cap \text{boîte réelle})}{\text{aire}(\text{boîte prédite} \cup \text{boîte réelle})} \in [0,1].
\end{equation}
Une IoU proche de $1$ indique une boîte prédite quasiment identique à la boîte
réelle ; une IoU proche de $0$ indique une absence de recouvrement. Cette mesure sert
à la fois à évaluer la qualité de la détection et, lors de l'entraînement, à
déterminer quelle prédiction est associée à quel objet réel.

Pour chaque boîte prédite, le réseau produit également un score de confiance
$\hat{c} \in [0,1]$, qui approche le produit
\begin{equation}
    \hat{c} \approx \Pr(\text{objet présent dans la cellule}) \times \text{IoU}(\text{boîte prédite}, \text{boîte réelle}).
\end{equation}
Ainsi, une détection n'est retenue que si un objet est jugé présent \emph{et} que la
boîte prédite recouvre bien l'objet réel.

\paragraph{Filtrage des détections redondantes.}

Plusieurs cellules voisines peuvent proposer des boîtes qui se recouvrent fortement
pour un même objet. Un post-traitement, la \textbf{suppression non-maximale} (NMS),
élimine les doublons : parmi les boîtes qui se recouvrent fortement (IoU élevé entre
elles), seule celle au score de confiance le plus élevé est conservée. C'est ce qui
garantit qu'un même visage n'est détecté qu'une seule fois dans l'image finale.

\paragraph{Paramétrisation des boîtes par rapport aux ancres.}

YOLO ne prédit pas directement les coordonnées d'une boîte, mais un ajustement par
rapport à une boîte de référence prédéfinie (\emph{anchor box}), de largeur $p_w$ et
hauteur $p_h$, associée à la cellule de centre $(c_x, c_y)$. Le réseau prédit quatre
valeurs brutes $(t_x, t_y, t_w, t_h)$, transformées en coordonnées réelles par :
\begin{equation}
    b_x = \sigma(t_x) + c_x, \qquad b_y = \sigma(t_y) + c_y,
\end{equation}
\begin{equation}
    b_w = p_w \, e^{t_w}, \qquad b_h = p_h \, e^{t_h},
\end{equation}
où $\sigma$ désigne la fonction sigmoïde, qui contraint le centre prédit à rester
dans la cellule considérée, tandis que l'exponentielle garantit la positivité des
largeur et hauteur prédites.

\paragraph{Fonction de perte composite, en détail.}

En notant $\mathbb{1}_{ij}^{\text{obj}}$ l'indicatrice valant $1$ si la cellule
$(i,j)$ contient effectivement le centre d'un objet réel (et $0$ sinon), et
$\mathbb{1}_{ij}^{\text{noobj}} = 1 - \mathbb{1}_{ij}^{\text{obj}}$, la perte totale
se décompose en trois termes :
\begin{equation}
    \mathcal{L}_{\text{YOLO}}(\theta) = \mathcal{L}_{\text{loc}}(\theta) + \mathcal{L}_{\text{obj}}(\theta) + \mathcal{L}_{\text{cls}}(\theta),
\end{equation}

\begin{equation}
    \mathcal{L}_{\text{loc}} = \lambda_{\text{coord}} \sum_{i,j} \mathbb{1}_{ij}^{\text{obj}} \Big[(b_x - \hat{b}_x)^2 + (b_y - \hat{b}_y)^2 + (\sqrt{b_w} - \sqrt{\hat{b}_w})^2 + (\sqrt{b_h} - \sqrt{\hat{b}_h})^2\Big],
\end{equation}

\begin{equation}
    \mathcal{L}_{\text{obj}} = \sum_{i,j} \mathbb{1}_{ij}^{\text{obj}} (\hat{c}_{ij} - c_{ij})^2 \; + \; \lambda_{\text{noobj}} \sum_{i,j} \mathbb{1}_{ij}^{\text{noobj}} \, \hat{c}_{ij}^{\,2},
\end{equation}

\begin{equation}
    \mathcal{L}_{\text{cls}} = \sum_{i,j} \mathbb{1}_{ij}^{\text{obj}} \sum_{k=1}^{K} -\log \hat{p}_{ij,k}.
\end{equation}

où $\mathcal{L}_{\text{loc}}$ pénalise l'erreur de régression sur les coordonnées de
la boîte englobante, $\mathcal{L}_{\text{obj}}$ pénalise l'erreur sur le score de
confiance $\hat{c}$ défini précédemment, et $\mathcal{L}_{\text{cls}}$ retrouve
exactement le terme d'entropie croisée établi en 1.1, restreint ici aux cellules
contenant réellement un objet.

Le passage par la racine carrée sur la largeur et la hauteur dans
$\mathcal{L}_{\text{loc}}$ est un choix technique : il pénalise proportionnellement
moins une erreur donnée sur une grande boîte que sur une petite boîte (une erreur de
10 pixels est plus grave, relativement, pour une petite boîte que pour une grande).
Le poids $\lambda_{\text{noobj}} < 1$ réduit l'importance des cellules sans objet
dans la perte, très majoritaires dans une image typique, pour éviter qu'elles ne
dominent l'apprentissage face aux rares cellules contenant effectivement un visage.

L'interprétation en termes de maximum de vraisemblance établie en 1.1 reste
valable, mais concerne uniquement la composante $\mathcal{L}_{\text{cls}}$ de cette
perte composite : c'est elle qui correspond à l'estimation de l'identité de
l'individu, tandis que $\mathcal{L}_{\text{loc}}$ et $\mathcal{L}_{\text{obj}}$
relèvent respectivement d'un problème de régression aux moindres carrés et d'un
problème de détection binaire (présence/absence d'objet), non directement concernés
par la classification des identités.

\emph{Note : les versions plus récentes de YOLO (v3 et suivantes) utilisent des
variantes plus sophistiquées, notamment pour $\mathcal{L}_{\text{loc}}$ (perte de
type CIoU) ; la formulation ci-dessus correspond à la version originale, suffisante
pour comprendre le principe général de la méthode.}

\section{Évaluation des performances du modèle}
\label{sec:eval-theorie}

Cette section présente, sur le plan théorique, les outils statistiques permettant
d'évaluer la qualité des prédictions d'un modèle de classification. Leur
application aux données réelles du stage est présentée au
chapitre~\ref{chap:proximite}, section~\ref{sec:eval-application}.

\paragraph{Matrice de confusion.}

Pour un problème de classification à $K$ classes, la \textbf{matrice de confusion}
$C$ est une matrice $K \times K$ dont l'entrée $C_{kl}$ compte le nombre
d'observations dont l'identité réelle est $k$ et l'identité prédite est $l$. La
diagonale de $C$ regroupe les prédictions correctes ; les termes hors diagonale
correspondent aux erreurs, et leur répartition renseigne sur les confusions
systématiques entre certains individus.

\paragraph{Métriques dérivées.}

À partir de la matrice de confusion, on définit pour chaque individu $k$ :
\begin{align}
    \text{précision}_k &= \frac{C_{kk}}{\sum_{l} C_{lk}} \quad \text{(parmi les prédictions "$k$", combien sont correctes)},\\
    \text{rappel}_k &= \frac{C_{kk}}{\sum_{l} C_{kl}} \quad \text{(parmi les individus réellement "$k$", combien sont bien retrouvés)},
\end{align}
ainsi que le taux d'exactitude global (\emph{accuracy}) :
\begin{equation}
    \text{exactitude} = \frac{\sum_k C_{kk}}{\sum_{k,l} C_{kl}}.
\end{equation}

\section{Utilisation de YOLO dans le cadre du projet}

La problématique posée -- détecter des macaques dans une vidéo, puis en déduire une
mesure de similarité/proximité entre individus -- requiert précisément les deux
capacités offertes conjointement par YOLO : la \textbf{localisation} (indispensable
pour extraire les coordonnées nécessaires au calcul ultérieur de la proximité
spatiale, cf. chapitre~\ref{chap:proximite}) et la \textbf{classification} (nécessaire
pour savoir quel individu est localisé à quel endroit). Le traitement en une seule
passe du réseau permet par ailleurs un traitement de vidéos entières, comportant
plusieurs dizaines de milliers de frames, à un coût de calcul raisonnable -- un
critère déterminant compte tenu du volume de données vidéo disponible sur la
population de Koshima (cf. chapitre~\ref{chap:donnees}). Le choix de ce modèle,
opéré par l'encadrant du stage en amont de celui-ci sur la base de résultats déjà
obtenus sur la population de Shodoshima, apparaît ainsi pleinement justifié au
regard de la problématique.


% =====================================================================
% CHAPITRE 2 — DONNÉES ET PRÉTRAITEMENT
% =====================================================================
\chapter{Données et prétraitement}
\label{chap:donnees}

\section{Présentation et acquisition des données}
% TODO : les vidéos ont été enregistrées sur le terrain (Koshima, Japon) par
% l'encadrant du stage ; expliquer pourquoi ces vidéos sont la seule source
% d'information permettant de mesurer la proximité sociale des individus, et
% en quoi leur exploitation par YOLO est nécessaire pour produire des données
% exploitables scientifiquement (cf. objectif AI-SARU, introduction)

\section{Organisation et préparation des données}

Les prédictions du modèle sont stockées au format JSON, selon une structure proche
du format COCO, organisée en trois blocs :
\begin{itemize}
    \item \texttt{images} : une entrée par frame vidéo (identifiant, nom de fichier,
    dimensions) ;
    \item \texttt{annotations} : une entrée par détection, contenant les coordonnées
    de la boîte englobante (\texttt{bbox}), l'identifiant de \emph{track}
    (\texttt{track\_id}), le score de confiance de la meilleure prédiction
    (\texttt{confidence}), l'identité retenue (\texttt{category\_id},
    \texttt{category\_name}), ainsi que le détail des cinq identités les plus
    probables et leurs scores respectifs (\texttt{extended\_classification},
    top1 à top5) ;
    \item \texttt{categories} : la liste des individus connus de la population
    étudiée (43 individus recensés dans l'exemple ci-dessous).
\end{itemize}

Un fichier JSON de ce type est généré par vidéo traitée. À titre d'illustration, un
fichier correspondant à une vidéo de la population de Koshima contient :
\begin{itemize}
    \item 68\,332 frames et 46\,647 détections (soit en moyenne 0,68 détection par
    frame, la plupart des frames ne contenant pas ou un seul individu détecté) ;
    \item 1\,270 \emph{tracks} distincts, de taille très variable (médiane : 11
    frames, maximum : 1\,171 frames) ;
    \item une confiance moyenne de la meilleure prédiction de 0,750
    (écart-type : 0,092) ;
    \item une forte hétérogénéité du nombre de détections par individu, allant de
    26 (individu le moins représenté) à 2\,670 (individu le plus représenté) ;
    \item une proportion notable de détections ambiguës : dans 29,6\% des cas,
    l'écart de score entre la première et la deuxième identité proposée est
    inférieur à 0,1, signe d'une identification incertaine du modèle.
\end{itemize}

\emph{Remarque méthodologique} : ce fichier ne contient pas de champ distinguant
explicitement une identité \emph{corrigée manuellement} d'une identité simplement
\emph{prédite par le modèle} -- \texttt{category\_name} correspond dans tous les cas
observés à la première proposition (\emph{top1}) de \texttt{extended\_classification}.
Une clarification reste nécessaire pour déterminer si les fichiers disponibles
correspondent aux prédictions avant ou après le ré-entraînement du modèle sur les
données corrigées de Koshima ; cette question est discutée en
section~\ref{sec:limites}.

Les vidéos sont par ailleurs numérisées en une séquence de frames à intervalle
régulier (échantillonnage temporel), chaque frame étant redimensionnée à une
résolution fixe avant d'être présentée en entrée du réseau.
% TODO : compléter avec les paramètres exacts (fréquence d'échantillonnage,
% résolution, normalisation des pixels) une fois confirmés

\section{Nettoyage et annotation des images}
% TODO :
% - protocole d'annotation et de correction manuelle via l'interface
% - critères de nettoyage (images floues, individus non identifiables même
%   manuellement, détections dupliquées, frames sans détection exploitable)

\section{Correction et validation des détections}

Afin d'évaluer objectivement la qualité des identifications produites par le
modèle, un échantillon d'images est relu manuellement : pour chaque image de
l'échantillon, l'identité réelle de l'individu est reportée dans un tableau, en
regard de l'identité prédite par le modèle. Ce tableau constitue la \textbf{vérité
terrain} nécessaire au calcul des métriques d'évaluation présentées au
chapitre~\ref{chap:methode}, section~\ref{sec:eval-theorie}, et appliquées au
chapitre~\ref{chap:proximite}.
% TODO : décrire concrètement la construction de ce tableau (nombre d'images
% relues, méthode d'échantillonnage, répartition entre les deux annotateurs)


% =====================================================================
% CHAPITRE 3 — CONSTRUCTION DES MATRICES DE PROXIMITÉ SOCIALE
% =====================================================================
\chapter{Construction des matrices de proximité sociale}
\label{chap:proximite}

\section{Détection des individus et évaluation des résultats}
\label{sec:eval-application}

\paragraph{Consolidation de l'identité par \emph{track.}}

Les détections ne sont pas traitées indépendamment image par image : elles sont
regroupées en \emph{tracks}, c'est-à-dire des séquences de détections correspondant à
un même individu suivi sur plusieurs frames consécutives. Pour un track $t$, on cumule
sur toutes ses frames les scores de confiance associés à chaque identité candidate $k$ :
\begin{equation}
    S_t(k) = \sum_{\text{frame } f \,\in\, t} \text{score}_f(k).
\end{equation}
L'identité retenue pour le track est celle qui maximise ce score cumulé, sous deux
conditions de fiabilité :
\begin{itemize}
    \item un support relatif suffisant, $\dfrac{\max_k S_t(k)}{\sum_k S_t(k)} \geq \tau_{\text{support}}$
    (seuil par défaut : $60\%$) ;
    \item un nombre minimal de frames dans le track (par défaut : 3).
\end{itemize}
Cette consolidation par vote pondéré agrégé rend l'identification plus robuste qu'une
décision frame par frame, chaque frame individuelle pouvant être bruitée.

\paragraph{Résultats d'évaluation.}
% TODO : appliquer les formules du chapitre 1 (matrice de confusion, précision,
% rappel, exactitude) aux données du tableau de vérité terrain construit au
% chapitre 2 ; présenter et critiquer les valeurs obtenues (bonnes ou mauvaises),
% comparer si possible avant/après ré-entraînement

\section{Extraction des coordonnées et calcul des distances}

Pour deux individus $a$ et $b$ détectés dans une même frame, on note
$(c_{x,a}, c_{y,a})$ et $(c_{x,b}, c_{y,b})$ les centres de leurs boîtes de détection
respectives. La distance centre-centre s'écrit
\begin{equation}
    d_{ab} = \sqrt{(c_{x,a}-c_{x,b})^2 + (c_{y,a}-c_{y,b})^2}.
\end{equation}

Cette distance en pixels ne suffit pas à elle seule, car elle ne tient pas compte de la
profondeur : un même écart en pixels peut correspondre à une grande distance réelle si
les individus sont loin de la caméra, ou une faible distance s'ils en sont proches. La
hauteur de la boîte de détection sert de proxy à l'échelle apparente de l'individu,
corrigée d'un facteur $\alpha(\text{âge})$ (les individus jeunes, plus petits, sont
sous-corrigés sinon) :
\begin{equation}
    h^{\text{corr}} = h \times \alpha(\text{âge}).
\end{equation}
La distance normalisée est alors
\begin{equation}
    d_{ab}^{\text{norm}} = \frac{d_{ab}}{\max(h_a^{\text{corr}}, h_b^{\text{corr}})}.
\end{equation}

Un ratio de taille des aires corrigées sert de garde-fou contre les incompatibilités de
profondeur :
\begin{equation}
    r_{ab} = \frac{\max(\text{aire}_a^{\text{corr}}, \text{aire}_b^{\text{corr}})}{\min(\text{aire}_a^{\text{corr}}, \text{aire}_b^{\text{corr}})}.
\end{equation}
Si $r_{ab}$ dépasse un seuil $\tau_{\text{ratio}}$, la paire est jugée trop incertaine
en profondeur et exclue.

La paire $(a,b)$ est déclarée \textbf{proche} dans la frame si
\begin{equation}
    d_{ab}^{\text{norm}} \leq \tau_{\text{dist}} \quad \text{et} \quad r_{ab} \leq \tau_{\text{ratio}}.
\end{equation}

\section{Construction de la matrice de proximité sociale à partir des détections}

\paragraph{Agrégation frame vers vidéo.}

Pour chaque vidéo $v$ et chaque paire $(a,b)$ co-détectée, on calcule la proportion de
frames où la paire est proche :
\begin{equation}
    \text{score}_v(a,b) = \frac{\text{nombre de frames "proches"}}{\text{nombre de frames co-détectées valides}}.
\end{equation}

\paragraph{Agrégation vidéo vers matrice globale.}

Les frames d'une même vidéo ne constituent pas des observations indépendantes : deux
individus filmés proches pendant 30 secondes ne fournissent pas autant d'information
que 900 observations indépendantes de proximité, mais plutôt une unique situation
répétée à haute fréquence d'échantillonnage. Pour limiter le poids disproportionné des
vidéos longues, chaque vidéo se voit attribuer un poids plafonné :
\begin{equation}
    w_v = \min(\text{nombre de frames co-détectées valides}_v, \; w_{\max}),
\end{equation}
avec $w_{\max} = 900$ par défaut (soit 30 secondes à 30 images/seconde).

La matrice de proximité globale $\hat{M}$ est alors une moyenne pondérée des scores
vidéo :
\begin{equation}
    \hat{M}_{ab} = \frac{\displaystyle\sum_{v} w_v \cdot \text{score}_v(a,b)}{\displaystyle\sum_{v} w_v}.
\end{equation}

\paragraph{Propriétés de la matrice obtenue.}
% TODO : symétrie, valeurs manquantes (paires jamais co-détectées), diagonale, échelle [0,1]

\section{Construction de la matrice de référence à partir des observations de terrain}

La matrice de référence $M^*$ est construite à partir d'observations de terrain
réalisées directement par des chercheurs, sur la même période temporelle que les
vidéos utilisées pour le ré-entraînement du modèle. Contrairement à la matrice
estimée $\hat{M}$, qui repose uniquement sur la fréquence de proximité spatiale
détectée automatiquement, $M^*$ agrège plusieurs dimensions du lien social observées
sur le terrain : la fréquence de coprésence, les interactions de toilettage
(\emph{grooming}), les relations de dominance, ainsi que les liens de parenté entre
individus. $M^*_{ab}$ traduit ainsi une intensité globale de proximité sociale entre
les individus $a$ et $b$, au sens large.

Cette différence de nature entre les deux matrices doit être nuancée. La matrice
estimée $\hat{M}$ ne mesure explicitement qu'une proximité spatiale ; elle ne
distingue pas la nature du lien à l'origine de cette proximité. Cependant, les
épisodes de grooming ou de jeu, qui impliquent par nature un contact ou une grande
proximité physique entre les individus concernés, se traduisent directement par des
scores de proximité élevés dans $\hat{M}$ : ces comportements sont donc captés
\emph{indirectement} par la mesure spatiale, sans être reconnus ou différenciés en
tant que tels. $\hat{M}$ ne peut en revanche pas distinguer une proximité spatiale
« fortuite » (par exemple deux individus qui se nourrissent côte à côte sans
interaction sociale particulière) d'une proximité liée à un comportement social
significatif, ni capter des liens sociaux ne se traduisant pas par une coprésence
spatiale fréquente (par exemple un lien de parenté entre deux individus qui se
croisent rarement). La matrice de référence $M^*$, construite à partir
d'observations de terrain intégrant explicitement la nature des interactions
(grooming, dominance, parenté), permet cette distinction que $\hat{M}$ ne fait pas.
Cette asymétrie doit être gardée à l'esprit dans l'interprétation de leur
comparaison (cf. section~\ref{sec:limites}).


% =====================================================================
% CHAPITRE 4 — COMPARAISON STATISTIQUE DES MATRICES
% =====================================================================
\chapter{Comparaison statistique des matrices}
\label{chap:comparaison}

\section{Méthodologie de comparaison des matrices}

Les entrées de la matrice $\hat{M}$ (et de la matrice de référence $M^*$) ne sont pas
indépendantes entre elles : la ligne associée à un individu $a$ particulièrement
sociable sera globalement élevée sur toute la ligne, car les valeurs $\hat{M}_{ab}$ et
$\hat{M}_{ac}$ partagent la même dépendance à $a$. Un test de corrélation classique
(Pearson), qui suppose l'indépendance des observations, ne peut donc pas être appliqué
directement entre les entrées des deux matrices. Une méthode adaptée à cette
structure de dépendance est nécessaire : le test de Mantel.

\section{Présentation théorique du test de Mantel}

Le \textbf{test de Mantel} \citep{mantel1967} est spécifiquement conçu pour comparer
deux matrices de distance ou de proximité définies sur le même ensemble d'unités
(ici, les $K$ individus).

\paragraph{Statistique de test.}

On considère la corrélation (par exemple de Pearson) entre les entrées des deux
matrices, restreinte aux paires distinctes $(a,b)$, $a \neq b$ :
\begin{equation}
    Z_{\text{obs}} = \text{corr}\big(\{\hat{M}_{ab}\}_{a<b}, \; \{M^*_{ab}\}_{a<b}\big).
\end{equation}

\paragraph{Hypothèse nulle et test par permutation.}

L'hypothèse nulle $H_0$ postule l'absence d'association entre les deux structures de
proximité. Comme les entrées d'une même matrice ne sont pas indépendantes (une ligne
entière dépend de la sociabilité globale d'un individu), la distribution théorique
usuelle de la statistique de corrélation ne s'applique pas. On construit donc la
distribution de $Z$ sous $H_0$ par \textbf{permutation} :
\begin{enumerate}
    \item on permute aléatoirement l'ordre des individus (lignes et colonnes
    simultanément) dans l'une des deux matrices, par exemple $M^*$, donnant $M^*_\pi$
    pour une permutation $\pi$ ;
    \item on recalcule la statistique de corrélation $Z_\pi$ entre $\hat{M}$ et
    $M^*_\pi$ ;
    \item on répète l'opération un grand nombre de fois (typiquement $999$ ou
    $9999$ permutations), donnant une distribution empirique $\{Z_\pi\}$ sous $H_0$.
\end{enumerate}

La p-valeur empirique est alors la proportion de permutations pour lesquelles
$Z_\pi$ est au moins aussi extrême que $Z_{\text{obs}}$ :
\begin{equation}
    p = \frac{1 + \#\{\pi : |Z_\pi| \geq |Z_{\text{obs}}|\}}{1 + \text{nombre de permutations}}.
\end{equation}

\section{Application du test de Mantel et résultats}

\paragraph{Vérification des conditions d'application.}

\begin{itemize}
    \item Les deux matrices doivent être définies sur le même ensemble d'individus,
    avec le même ordre de lignes/colonnes avant permutation.
    \item Le nombre de permutations doit être suffisant pour estimer précisément la
    p-valeur (au minimum quelques centaines).
    \item Les valeurs manquantes (paires jamais co-détectées dans $\hat{M}$) doivent
    être traitées explicitement (exclusion de la paire dans les deux matrices, ou
    imputation à justifier).
    \item Le test de Mantel suppose que la structure de dépendance interne à chaque
    matrice est correctement neutralisée par la permutation conjointe des
    lignes/colonnes -- ce qui est le cas ici, la permutation appliquée à un individu
    déplace simultanément toute sa ligne et sa colonne.
\end{itemize}

\paragraph{Résultats obtenus.}
% TODO : coefficient de corrélation observé, p-valeur, nombre de permutations
% utilisées, individus/paires exclus le cas échéant

\section{Interprétation des résultats}
% TODO : le résultat du test permet-il de conclure à une association entre les
% deux matrices ? Discuter en tenant compte de la nuance établie au chapitre 3
% (proximité spatiale vs. lien social composite)


% =====================================================================
% CHAPITRE 5 — DISCUSSION ET PERSPECTIVES
% =====================================================================
\chapter{Discussion et perspectives}

\section{Discussion des résultats}
% TODO : synthèse croisée -- qualité de détection (chapitre 3.1) et fidélité de
% la matrice de proximité (chapitre 4) ; propagation des erreurs d'identification
% vers la matrice de proximité ; apport personnel et degré d'autonomie pendant
% le stage

\section{Limites et conditions d'application des méthodes}
\label{sec:limites}
% TODO :
% - absence d'ensemble de test indépendant (train/validation seulement)
% - non-indépendance temporelle des frames au sein d'une vidéo
% - hypothèses du proxy de profondeur (ratio de taille des boîtes)
% - taille de l'échantillon Koshima
% - incertitude sur le statut avant/après ré-entraînement des fichiers disponibles
% - échantillon de vérité terrain limité (chapitre 2)

\section{Perspectives d'amélioration}
% TODO : poursuite du ré-entraînement avec davantage de données, extension à
% d'autres populations, amélioration du proxy de profondeur, automatisation de
% la constitution de la vérité terrain, etc.


% =====================================================================
% CONCLUSION GÉNÉRALE
% =====================================================================
\chapter*{Conclusion générale}
\addcontentsline{toc}{chapter}{Conclusion générale}
% TODO : synthèse des principaux résultats, retour sur la problématique posée
% en introduction


% ===================== BIBLIOGRAPHIE =====================
\addcontentsline{toc}{chapter}{Bibliographie}
\bibliographystyle{plainnat}
\bibliography{biblio}


% ===================== ANNEXES =====================
\appendix
\chapter{Annexes}

\section{Extraits de code}
% Rappel des consignes : lignes principales uniquement, jamais le programme complet
\begin{lstlisting}
def evaluate_pair_in_frame(a, b, dist_thresh_norm, size_ratio_thresh):
    dist_norm = distance(a, b) / max(a["corrected_h"], b["corrected_h"])
    size_ratio = max(a["aire"], b["aire"]) / min(a["aire"], b["aire"])
    close = (dist_norm <= dist_thresh_norm) and (size_ratio <= size_ratio_thresh)
    return close
\end{lstlisting}

\section{Figures complémentaires}
% TODO : heatmap de la matrice, visualisation du réseau social

\section{Tableaux complémentaires}
% TODO : matrice de confusion complète (chapitre 3)

\end{document}